% GENERATED FILE. Edit canonical lessons, then run npm run generate:books.
% canonical-source-hash: fnv1a64:cbbabf73e1a0ef7e
% canonical-lessons: ZH-C75-kuaile, ZH-W75-le-joy, ZH-C75-shengri, ZH-C75-dangao, ZH-C75-liwu, ZH-C75-song

\chapter{Celebrating}
\label{ch:zh-celebrating}
\hlchaptermodality{75}

\begin{chapteropening}
\textbf{By the end of this chapter:} \emph{I can wish somebody a happy birthday and give a present.}

Complete the last lesson of chapter 75: I can wish somebody a happy birthday and give a present.
\end{chapteropening}

\section[\texorpdfstring{\zh{快乐} (kuàilè)}{快乐 (kuàilè)}]{\zh{快乐} (kuàilè) — joyful; happy}

\label{lesson:ZH-C75-kuaile}



\begin{quote}
Before the new one: say the Mandarin for welcome, then the Mandarin for to introduce.
\end{quote}

\subsection*{You'll want to know: \zh{快乐}}
\textbf{\zh{快乐}} — \emph{kuàilè} — \textquotedblleft{}joyful; happy\textquotedblright{}.

\begin{quote}
\textbf{\zh{我们很快乐}} \emph{wǒmen hěn kuàilè} — We're very happy.
\end{quote}

The new character, in pieces: \zh{乐} — a nearly flat fall, a turn, and then \zh{小} with its middle hooked. The next lesson puts it in your hand.

\begin{cousinweb}
\textbf{\zh{快乐}} (\emph{kuàilè}) is built from \zh{快} and \zh{乐}, both inherited from Middle Chinese, \emph{khwaejH} and \emph{lak} (Guangyun, in Baxter's transcription). The dictionary glosses the first as 'glad' and the second as 'joy': gladness twice over.
\end{cousinweb}

\begin{grammarlens}[title={What you've built}]
One more word for meeting people.
\end{grammarlens}

\subsection*{Your turn}
\begin{itemize}
\raggedright
  \item \emph{Say it:} \emph{kuàilè}
  \item \emph{Say it:} \emph{kuàilè}, once more
  \item \emph{Say it:} \emph{wǒmen hěn kuàilè} — We're very happy.
  \item \emph{Recall:} say the Mandarin for welcome, then the Mandarin for to introduce, then say \emph{kuàilè} again
\end{itemize}

\begin{tcolorbox}[breakable,colback=teal!4,colframe=teal!35!black,title={Before you move on}]
What does \zh{快乐} mean? (\textquotedblleft{}Joyful; happy\textquotedblright{}.) Say it once more.
\end{tcolorbox}
\section[\texorpdfstring{\zh{乐} (lè)}{乐 (lè)}]{\zh{乐} — a fall, a turn, and \zh{小}}

\label{lesson:ZH-W75-le-joy}



\begin{quote}
The word first, then the character inside it.

\begin{itemize}
\raggedright
  \item \emph{Recall:} say \emph{kuàilè} — \textbf{R1}, one lesson back, before it decays
\end{itemize}

\begin{quote}
\zh{小} \zh{快}
\end{quote}

\emph{Read it:} each one aloud, with nothing said first — \emph{xiǎo}, \emph{kuài}
\end{quote}

\subsection*{Script — a fall, a turn, and \zh{小}}
\hlblockfigure{\detokenize{figures/ZH-W75-le-joy-filmstrip.pdf}}{How \zh{乐} is written, stroke by stroke}

\begin{quote}
\zh{乐}
\end{quote}

\textbf{One piece, five strokes.}

\textbf{A left-falling stroke} across the top, from right to left, nearly flat.

\textbf{A vertical that turns right}, under the left end of the first stroke.

\textbf{Then the three strokes of \zh{小}}: the vertical down the middle that hooks to the left, a dot falling left, a dot falling right.

2 + 3 = five.

\subsection*{Your turn}
\begin{itemize}
\raggedright
  \item \emph{Write it:} \zh{乐}, all five, with the model in front of you
  \item \emph{Cover:} write it again from memory, and count the pieces, not the strokes
\end{itemize}

\begin{tcolorbox}[breakable,colback=teal!4,colframe=teal!35!black,title={Before you move on}]
How many strokes in \zh{乐}? (\textbf{Five}.) Which character you know are its last three? (\textbf{\zh{小}}.) Which way does the first stroke run? (\textbf{From right to left, nearly flat}.)

Source: \href{https://raw.githubusercontent.com/chanind/hanzi-writer-data/68d10a4b21150cae5e1ebbd223eed289cf32d90c/data/\%E4\%B9\%90.json}{Hanzi Writer Data \zh{乐}, pinned stroke paths and medians}.
\end{tcolorbox}
\section[\texorpdfstring{\zh{生日} (shēngrì)}{生日 (shēngrì)}]{\zh{生日} (shēngrì) — birthday}

\label{lesson:ZH-C75-shengri}



\begin{quote}
Before the new one: say the Mandarin for to introduce, then the Mandarin for joyful.
\end{quote}

\subsection*{You'll want to know: \zh{生日}}
\textbf{\zh{生日}} — \emph{shēngrì} — \textquotedblleft{}birthday\textquotedblright{}.

\begin{quote}
\textbf{\zh{明天是我的生日}} \emph{míngtiān shì wǒ de shēngrì} — Tomorrow is my birthday.
\end{quote}

\begin{cousinweb}
\textbf{\zh{生日}} (\emph{shēngrì}) is built from \zh{生} 'to be born' and \zh{日} 'day', both inherited from Middle Chinese, \emph{srjaeng} and \emph{nyit} (Guangyun, in Baxter's transcription).
\end{cousinweb}

\begin{grammarlens}[title={What you've built}]
One more word for meeting people.
\end{grammarlens}

\subsection*{Your turn}
\begin{itemize}
\raggedright
  \item \emph{Say it:} \emph{shēngrì}
  \item \emph{Say it:} \emph{shēngrì}, once more
  \item \emph{Say it:} \emph{míngtiān shì wǒ de shēngrì} — Tomorrow is my birthday.
  \item \emph{Recall:} say the Mandarin for to introduce, then the Mandarin for joyful, then say \emph{shēngrì} again
\end{itemize}

\begin{tcolorbox}[breakable,colback=teal!4,colframe=teal!35!black,title={Before you move on}]
What does \zh{生日} mean? (\textquotedblleft{}Birthday\textquotedblright{}.) Say it once more.
\end{tcolorbox}
\section[dàngāo]{dàngāo — cake}

\label{lesson:ZH-C75-dangao}



\begin{quote}
Before the new one: say the Mandarin for joyful, then the Mandarin for birthday.
\end{quote}

\subsection*{You'll want to know: dàngāo}
\emph{dàngāo} — \textquotedblleft{}cake\textquotedblright{}.

\begin{quote}
\emph{māma de dàngāo hěn hǎochī} — Mum's cake is delicious.
\end{quote}

Its characters come later; for now it is a word for the ear.

\begin{cousinweb}
\emph{dàngāo} is built from dàn 'egg', the dàn of jīdàn, and gāo 'cake'. No dictionary we use records an older form of the first, so its history is not settled; the second is probably the \emph{kaw} that the Guangyun writes with another form of the character (Guangyun, in Baxter's transcription), glossed with a word for gruel.
\end{cousinweb}

\begin{grammarlens}[title={What you've built}]
One more word for meeting people.
\end{grammarlens}

\subsection*{Your turn}
\begin{itemize}
\raggedright
  \item \emph{Say it:} \emph{dàngāo}
  \item \emph{Say it:} \emph{dàngāo}, once more
  \item \emph{Say it:} \emph{māma de dàngāo hěn hǎochī} — Mum's cake is delicious.
  \item \emph{Recall:} say the Mandarin for joyful, then the Mandarin for birthday, then say \emph{dàngāo} again
\end{itemize}

\begin{tcolorbox}[breakable,colback=teal!4,colframe=teal!35!black,title={Before you move on}]
What does dàngāo mean? (\textquotedblleft{}Cake\textquotedblright{}.) Say it once more.
\end{tcolorbox}
\section[lǐwù]{lǐwù — present; gift}

\label{lesson:ZH-C75-liwu}



\begin{quote}
Before the new one: say the Mandarin for birthday, then the Mandarin for cake.
\end{quote}

\subsection*{You'll want to know: lǐwù}
\emph{lǐwù} — \textquotedblleft{}present; gift\textquotedblright{}.

\begin{quote}
\emph{wǒ mǎi le yí ge lǐwù} — I've bought a present.
\end{quote}

Its characters come later; for now it is a word for the ear.

\begin{cousinweb}
\emph{lǐwù} is built from lǐ 'rite, courtesy' and wù 'thing', both inherited from Middle Chinese, \emph{lejX} and \emph{mjut} (Guangyun, in Baxter's transcription): a thing given out of courtesy.
\end{cousinweb}

\begin{grammarlens}[title={What you've built}]
One more word for meeting people.
\end{grammarlens}

\subsection*{Your turn}
\begin{itemize}
\raggedright
  \item \emph{Say it:} \emph{lǐwù}
  \item \emph{Say it:} \emph{lǐwù}, once more
  \item \emph{Say it:} \emph{wǒ mǎi le yí ge lǐwù} — I've bought a present.
  \item \emph{Recall:} say the Mandarin for birthday, then the Mandarin for cake, then say \emph{lǐwù} again
\end{itemize}

\begin{tcolorbox}[breakable,colback=teal!4,colframe=teal!35!black,title={Before you move on}]
What does lǐwù mean? (\textquotedblleft{}Present; gift\textquotedblright{}.) Say it once more.
\end{tcolorbox}
\section[sòng]{sòng — to give as a present; to see off}

\label{lesson:ZH-C75-song}



\begin{quote}
Before the new one: say the Mandarin for cake, then the Mandarin for present.
\end{quote}

\subsection*{You'll want to know: sòng}
\emph{sòng} — \textquotedblleft{}to give as a present; to see off\textquotedblright{}.

\begin{quote}
\emph{wǒmen sòng shénme lǐwù} — What present shall we give?
\end{quote}

Its characters come later; for now it is a word for the ear.

\begin{cousinweb}
\emph{sòng} is inherited from Middle Chinese \emph{suwngH} (Guangyun, in Baxter's transcription), glossed 'to send off'; giving a present is sending something off with somebody.
\end{cousinweb}

\begin{grammarlens}[title={What you've built}]
One more word for meeting people.
\end{grammarlens}

\subsection*{Your turn}
\begin{itemize}
\raggedright
  \item \emph{Say it:} \emph{sòng}
  \item \emph{Say it:} \emph{sòng}, once more
  \item \emph{Say it:} \emph{wǒmen sòng shénme lǐwù} — What present shall we give?
  \item \emph{Recall:} say the Mandarin for cake, then the Mandarin for present, then say \emph{sòng} again
\end{itemize}

\begin{tcolorbox}[breakable,colback=teal!4,colframe=teal!35!black,title={Before you move on}]
What does sòng mean? (\textquotedblleft{}To give as a present; to see off\textquotedblright{}.) Say it once more.
\end{tcolorbox}
